Pour fabriquer un électroaimant, on enroule \num{\v{N}} spires de fil de cuivre sur une longueur de \qty{\v{L}}{cm}. Le fil a une résistance totale de \qty{\v{R}}{\ohm} et on le branche sur une pile de \qty{\v{U}}{V}. La perméabilité du vide est $\mu_0 = \qty{1.257e-6}{T.m/A}$ (on néglige l'effet d'un noyau).

\begin{enumerate}[label=\alph*)]
    \item Déterminez l'intensité du courant dans la bobine.
    \item Déterminez le champ magnétique au centre de la bobine, en millitesla.
    \item Nommez deux moyens d'augmenter le champ magnétique de cet électroaimant sans changer la pile.
\end{enumerate}
